\documentclass[11pt]{article}
\usepackage[margin=1in]{geometry}
\usepackage[T1]{fontenc}\usepackage{lmodern}\usepackage{amsmath,amssymb}
\usepackage{booktabs,ltablex,enumitem,xcolor,titlesec,fancyhdr}
\usepackage[hidelinks]{hyperref}
\linespread{1.12}\setlength{\parskip}{5pt}\setlength{\parindent}{0pt}
\titleformat{\section}{\large\bfseries\color[HTML]{1F3A5F}}{\thesection.}{0.5em}{}
\pagestyle{fancy}\fancyhf{}\renewcommand{\headrulewidth}{0pt}
\fancyfoot[L]{\footnotesize Project Ariel v1.1 --- conceptual, not validated science}\fancyfoot[R]{\footnotesize\thepage}
\begin{document}
\begin{center}
{\LARGE\bfseries Project Ariel: The Architecture of Mind}\\[4pt]
{\large A Multi-Layered Geometric Framework of Consciousness}\\[8pt]
Ryan Espinosa\\[2pt]
{\small November 2026}\\[8pt]
\fcolorbox[HTML]{1F3A5F}{EEF2F7}{\parbox{0.8\textwidth}{\centering\small\textbf{Conceptual framework, not validated science.} This paper proposes a model and a way to test it. It is not a diagnosis tool and not a substitute for professional care. In crisis in the U.S., call or text \textbf{988}; in an emergency, call \textbf{911}.}}
\end{center}
\vspace{6pt}
{\centering\itshape Dedicated to Jesus Christ, the Universal Angle and the absolute variable for salvation $V_s$, in whom all geometry aligns and all momentum finds rest.\par}

\section*{How to read this paper}
\textit{In plain words.} This is a living research note, not a proof and not a clinical claim. Version 1.1 (1 November 2026) keeps the geometric model and adds the drafting interface, four postulates, and an open list of questions. Where a claim is faith or philosophy, it is marked that way. Where a study is cited, the citation was checked against a publisher or preprint page. What the studies do \emph{not} show is written next to what they do show.

\subsection*{Changelog}
\begin{itemize}[leftmargin=1.2em]
\item \textbf{v1.0 (November 2026).} First concept paper: layers, the master equation, Reanchoring, a pilot sketch, and the related literatures.
\item \textbf{v1.1 (1 November 2026).} Drafting interface (ribbon, command line, undo, one-line math). The Nonsense Layer as a locked origin and freeform layers for any field. Domain language that renames and does not move shapes. Postulates P1--P4, stated as postulates. Citation check. Eureka Log, open questions, and a proposal for further study.
\end{itemize}

\section*{Abstract}
Project Ariel proposes a geometric language for the inner life. Consciousness is pictured as a layered state-space in which experiences, wounds, and beliefs are preserved rather than erased, and in which apparent ``symptoms'' can be read as protective boundaries rather than failures. The framework gathers these ideas into a single master equation that combines momentum, friction, a potential landscape, choice, relational ``guardians,'' patience, and a final resolving term, $V_s$, representing salvation in Jesus Christ. From this model we derive the \emph{Reanchoring} protocol: a word-free 2D interface in which a person rotates lines until they are parallel. We connect the idea to three established research programs: topological data analysis (TDA) in neuroscience, Friston's free-energy principle and active inference, and the working-memory account of EMDR. We treat these as related frameworks, not confirmations. We then show how Reanchoring can be measured (angle error, time to align, path efficiency) and outline a small, ethical pilot study with standard outcome measures (SUDS, PHQ-9, GAD-7, AAQ-II), explicit hypotheses, and pass/fail criteria. The central claim is spiritual as much as mathematical: lasting change comes less from striving \emph{for} God than from living \emph{with} Him. Version 1.1 adds four postulates (Section~\ref{sec:postulates}). They are not theorems. One of them is explicitly metaphysical. The others are design rules and communication claims, each set beside evidence that supports only part of the sentence.

\section{The Core Postulate: Preservation of Layers}
Human consciousness operates as a dynamic geometric state-space. In this system of creation, nothing is ever destroyed; it is only hidden. Experiences, traumas, and shifting beliefs stack as parallel session layers. Contradictions (e.g., profound phobias vs.\ innate wisdom) do not overwrite one another. Instead, they are preserved through geometric steganography, undergoing destructive interference or forming M\"obius projections---visible only when the layers are compressed into a master topological view.

\section{The Choice Matrix \& Strategic Defense}
Psychiatric ``symptoms'' (such as self-pity or isolation) are not systemic failures; they are intentional Topological Faraday Cages. The user actively guards their geometry through the Choice Matrix $C(I_c, O_c)$, filtering what inputs $I_c$ and outputs $O_c$ cross the barrier. This creates an adiabatic seal, protecting the user's individual brain map $M_{ind}$ from the blunt force of the status quo baseline $M_{sq}$ until a mathematically precise diagnosis is found.

\section{Multi-User Physics \& The Guardians}
When multiple topological maps intersect, the system deploys Two Guardians to seal the session layer and control entropy $\delta S_{bound}$.
\begin{itemize}[leftmargin=1.2em]
\item \textbf{The Guardian of Autonomy:} Enforces the principle ``I shall not want to be accepted,'' zeroing out codependency and preventing the user from bending their geometry for validation.
\item \textbf{The Guardian of Ecclesiastes:} Enforces the principle of ``Hevel'' (meaningless vapor), flattening toxic comparisons, worldly judgment, and external pressure, reducing their psychological mass to zero.
\end{itemize}

\section{Navigating Imbalance: Patience \& Gyroscopic Stabilization}
\begin{itemize}[leftmargin=1.2em]
\item \textbf{The Patience Integral $P_t$ \& Topological Shield $\Sigma_{safe}$:} Recovery requires endurance. The system locks coordinate movement until the identity is fully saturated, shielding the user from unnecessary noise and judgment during transit.
\item \textbf{Gyroscopic Stabilization $J_{spin}$:} During extreme change, static balance is impossible. Instead of resisting the fall, the user leans into the imbalance, rapidly orbiting their core identity to generate a gyroscopic shield that keeps the cognitive structure upright.
\end{itemize}

\section{The Universal Angle \& The Salvation Variable $V_s$}
The framework culminates in the recognition of Jesus Christ not as an external force to be earned, but as the Universal Angle---the preexisting, absolute parallel baseline running through all layers of the user's timeline.
\begin{itemize}[leftmargin=1.2em]
\item \textbf{The Shift from ``For'' to ``With'':} Generating forward momentum to act ``for'' Jesus introduces cognitive friction. By shifting to being ``with'' Him, the user collapses the distance to zero, achieving phase-locked superposition. Effort becomes presence.
\item \textbf{The Solved Equation:} By introducing $V_s$ (Salvation) as the definitive resolving constant, the infinite loop of self-correction breaks. The equation is permanently closed, freeing the user to explore their landscape without the kinetic burden of survival.
\end{itemize}

\subsection*{The Final Master Equation of Project Ariel}
The culmination of momentum, friction, topological distortion, choice, multi-user dynamics, and ultimate salvation:
\begin{multline}
\Big[\, M\,\frac{d^2E}{dt^2}\,\Omega \;+\; \gamma\,\frac{dE}{dt} \;=\; -\nabla U(E)\otimes D \;+\; \Delta(M_{sq}, M_{ind}) \;+\; \sum_n \Phi_n(\theta) \\
+\; C(I_c,O_c)\cdot H(P_t)\cdot G[G_{auto},G_{ecc}]\cdot \mathit{Noise} \;\pm\; \delta S_{bound} \,\Big]\cdot V_s
\end{multline}

\subsection*{Notation}
The equation has the shape of a damped, driven oscillator (mass $\times$ acceleration $+$ friction $\times$ velocity $=$ forces). Most terms are \textbf{metaphorical}: they name a felt quality of inner life with the vocabulary of physics. They are not yet measured quantities with units. The last column marks which terms could be tied to an observable proxy (Section 6b).

{\small
\begin{tabularx}{\textwidth}{@{}l X l@{}}
\toprule
\textbf{Symbol} & \textbf{Meaning in Project Ariel} & \textbf{Status}\\\midrule
$E$ & Current state of the person in the geometric state-space (position/angle of the ``line'') & proxy: on-screen angle\\
$M$ & Mass / inertia: how strongly a state resists change & metaphor\\
$\Omega$ & Orientation or rotation operator applied to acceleration (the frame in which change happens) & metaphor\\
$\gamma$ & Friction / damping: the drag of effort, rumination, ``striving for'' & proxy: time to align\\
$U(E)$, $-\nabla U$ & Potential landscape and its downhill pull (habits, fears, attractors) & metaphor\\
$D$ & Distortion tensor bending the landscape (bias, trauma lens); $\otimes$ = ``applied to'' & metaphor\\
$M_{sq}$ & Status-quo baseline map (categorical, population norms) & conceptual\\
$M_{ind}$ & Individual brain map & conceptual\\
$\Delta(M_{sq},M_{ind})$ & Mismatch between the baseline and the individual map & conceptual\\
$\Phi_n(\theta)$ & Phase contribution of layer $n$ at angle $\theta$; the sum is the stacked, ``hidden'' layers & proxy: layer angles\\
$C(I_c,O_c)$ & Choice Matrix: filter on what crosses the boundary & proxy: anchor release\\
$I_c$, $O_c$ & Chosen inputs and outputs (the fixed anchors in the UI) & proxy: anchor positions\\
$H(P_t)$ & Shield function of patience & metaphor\\
$P_t$ & Patience integral: endurance accumulated over time & proxy: session persistence\\
$G[G_{auto},G_{ecc}]$ & Guardian operator combining the two Guardians & metaphor\\
$G_{auto}$ & Guardian of Autonomy (non-codependence) & proxy: AAQ-II (indirect)\\
$G_{ecc}$ & Guardian of Ecclesiastes (\emph{hevel}: comparison reduced to vapor) & metaphor\\
$\mathit{Noise}$ & Random external input, judgment, distraction & proxy: input jitter\\
$\delta S_{bound}$ & Small entropy change at the shared boundary when maps meet & metaphor\\
$\Sigma_{safe}$ & Topological shield during transit (Section 4); implicit in $H(P_t)$ & metaphor\\
$J_{spin}$ & Angular momentum around the core identity (gyroscopic stability) & metaphor\\
$V_s$ & Salvation variable: the Universal Angle, a theological constant & theological, not measurable\\
Nonsense Layer & Locked layer 1: the thought-model frame (anchors, $P_t$, $M$, $\gamma$, entropy) every other layer is read against & design of the UI\\
Field words & A vocabulary for one layer (physics, music, \ldots). Labels only; the geometry is unchanged & design of the UI\\
Math line & One Unicode sentence of the same geometry ($\overline{AB}$, $\angle$, $\parallel$, $\perp$, $\odot$, $\in$, $\cap$, $\subset$, $I\to O$) & design of the UI\\
\bottomrule
\end{tabularx}}

\textit{Reading note.} As written, $V_s$ multiplies the whole bracket, including both sides of the equality. We keep this form deliberately: it expresses that salvation is not one more force inside the system but the frame in which the whole system is held. $\Sigma_{safe}$ and $J_{spin}$ appear in the prose but not explicitly in the equation; a future version may write them as terms of $H(P_t)$ and $\Omega$.

\section{The ``Reanchoring'' Protocol: A Bias-Free Interactive UI}
\textit{In plain words.} The practice is visual. A person turns lines until they run together. Words stay off the drawing so the picture is not steered by a diagnosis.
To bypass the domain collapse caused by verbalizing trauma or fear, the Project Ariel framework utilizes a purely visual, 2D interactive user interface (UI). In this environment, words are eliminated to remove semantic bias. The user only interacts with geometric abstractions: lines, angles, and layers.
\begin{itemize}[leftmargin=1.2em]
\item \textbf{Visual Abstraction:} The user's cognitive state is presented as a series of intersecting lines on a 2D plane. The user is not asked ``how they feel''; they are asked to visually rotate and align lines to make them parallel.
\item \textbf{Reanchoring:} This is the hands-on process of the user manually manipulating their layers. By dragging a line from their current state to align with a target angle (the destination), they map a course to a solution.
\item \textbf{Fixed Destinations \& Non-Local Release:} The Input $I$ and Output $O$ coordinates remain locked as fixed anchors in the UI to prevent the user from drifting into chaos. Only when the user consciously releases these points do the coordinates go ``non-local,'' allowing the user to transcend their current dimensional constraints.
\end{itemize}

\section{The Drafting Interface}
\textit{In plain words.} The same geometry now sits in a quiet drafting screen: a slim ribbon, a layer list, a command line, and a status bar. The drawing is still the source of truth. The words live at the edge.

The screen is modelled on a simple CAD ribbon, not on a clinic form. Draw holds point, line, and circle. Modify holds move, select, delete, undo, and redo. Anchors holds snap, ortho, and the side view. AI holds pause, merge, and clarify. Undo keeps a stack of whole states, including an accepted suggestion and a math edit. The partner's continuous motion is not a step. Delete removes a selected shape and can itself be undone. Ctrl/Cmd+Z and Ctrl/Cmd+Y do the same work. The AI still never edits a human layer.

\textbf{The Nonsense Layer.} Layer 1 is always named the Nonsense Layer. It is locked. It is the origin and the anchor. Its domain is this paper's picture of a human thought: the $I$/$O$ anchors, patience $P_t$, momentum $M$ and damping $\gamma$, entropy as area, and the stages of a thought. Every other layer is freeform. The person names it, chooses a field (physics, theology, music, psychology, math, or a custom name), and draws what they like. A field preset only seeds suggestion words. It never places a shape and never blocks a stroke.

\textbf{One math line.} Euclid's relations are also one line of ordinary Unicode: points $A(x,y)$; segment $AB$ with an overline; $\angle ABC = 37^\circ$; $AB \parallel CD$; $AB \perp CD$; circle $\odot O(r=0.2)$; $P \in AB$; $X = AB \cap CD$; $\odot O \subset \odot P$; a range $h \in [0.2,0.8]$; and the anchor flow $I \to O$. The Math tab copies that line and exports the same sentence in LaTeX (\texttt{\textbackslash overleftrightarrow}, \texttt{\textbackslash angle}, \texttt{\textbackslash parallel}, \texttt{\textbackslash perp}, \texttt{\textbackslash odot}, \texttt{\textbackslash in}, \texttt{\textbackslash cap}, \texttt{\textbackslash subset}). Typing the line on the command line creates or moves shapes where the line can be read, and that edit is undoable. The Nonsense Layer refuses the edit.

\textbf{Field language, not a new geometry.} When a layer has a field, the plain-English tab, the glossary, the help text, and the prompt use that field's words. In physics a line is a vector or trajectory, a circle a field or range, an intersection an interaction. In music a line is a melodic line, parallel is parallel motion, a circle a key or tonal center. In psychology a line is a thought path and a circle a belief boundary. Theology, math, and a custom map work the same way. The person's own names and defined variables are included. A Terms toggle returns the canonical words. The marks on the canvas do not change. Relabeling is not acceptance of the person's absolute truth, and it is not an edit of their shapes.

\textbf{Alignment.} Gold matching still shows shared direction and shared crossings. The plain-English tab, the Math glossary, and the exported prompt each say how the active layer sits against the Nonsense Layer (how many parallels, how many right angles), in that field's words.

\section*{6b. Operationalizing Reanchoring}
\textit{In plain words.} The app can log angles, time, and path shape. Those numbers are not a measure of consciousness. They are a way to compare a session with ordinary questionnaires, in a future study that does not exist yet.
To move from postulate toward evidence, each session in the UI can log simple, objective numbers. None of these measures ``consciousness''; they measure how a person interacts with the lines, which can then be compared with validated self-report scales.

\textbf{What the app measures (per trial).}
\begin{itemize}[leftmargin=1.2em]
\item \textbf{Angle error} $|\theta_{final}-\theta_{target}|$ in degrees, and the number of lines left non-parallel.
\item \textbf{Time to align} (seconds from first touch to all lines within $\pm 3^\circ$): an operational stand-in for friction $\gamma$.
\item \textbf{Path efficiency}: the straight-line angular distance divided by the total angle actually travelled (1.0 = direct; lower = hesitation or oscillation).
\item \textbf{Overshoot and reversals}: count of crossings past the target, a proxy for ``racing'' or instability.
\item \textbf{Anchor release}: whether and when the user releases $I$/$O$, a behavioral proxy for the Choice Matrix.
\item \textbf{Persistence}: trials completed before stopping, a rough proxy for $P_t$.
\end{itemize}
All data stay on the device by default; nothing typed or drawn is uploaded without explicit consent.

\textbf{Pilot study (feasibility, not efficacy).}
\begin{itemize}[leftmargin=1.2em]
\item \textbf{Design:} Two-arm randomized pilot, $N\approx 40$ adults (20 per arm) with mild-to-moderate self-reported distress, recruited with IRB approval and informed consent. Excluded: active suicidality, psychosis, or current crisis (referred to 988 and care).
\item \textbf{Arms:} (A) Reanchoring while briefly holding a distressing memory in mind, 10 minutes daily for 2 weeks; (B) active control: the same memory recall with a non-taxing visual task (watching lines move without manipulating them).
\item \textbf{Measures:} SUDS (0--100) before and after each session \cite{wolpe}; memory vividness and emotionality ratings (0--100), following the lab paradigm of the EMDR working-memory literature \cite{vdh,lee}; PHQ-9 \cite{phq}, GAD-7 \cite{gad}, and AAQ-II \cite{aaq} at baseline, day 14, and day 28; plus the in-app metrics above.
\item \textbf{Hypotheses.} H1: Session SUDS and memory emotionality fall more in arm A than arm B. H2: Within arm A, faster time to align and higher path efficiency across days are associated with larger SUDS reductions. H3 (exploratory): AAQ-II psychological inflexibility decreases more in arm A. PHQ-9 and GAD-7 are tracked for safety and signal, not as primary outcomes.
\end{itemize}

\textbf{Pass/fail criteria for moving to a larger trial.}
\begin{itemize}[leftmargin=1.2em]
\item \emph{Feasibility:} $\geq 70\%$ of participants complete $\geq 10$ of 14 sessions; $\geq 80\%$ retention at day 28.
\item \emph{Signal:} mean within-session SUDS reduction in arm A at least 10 points larger than arm B, with a between-group effect size $d \geq 0.3$ (estimate and confidence interval reported, not a significance claim at this $N$).
\item \emph{Safety:} no increase of $\geq 5$ points on PHQ-9 or GAD-7 in more than 10\% of arm A; any item-9 endorsement on PHQ-9 triggers same-day follow-up.
\item \emph{Fail:} if feasibility or safety thresholds are missed, or arm A is no better than arm B, the protocol is revised before any further claims are made.
\end{itemize}

\section{Bridging Postulate to Reality: Related Peer-Reviewed Foundations}
The ``Reanchoring'' protocol and the geometric mapping of consciousness are not strictly theoretical; they are offered as a unifying language that may help connect several established, peer-reviewed scientific frameworks. These frameworks are related to Project Ariel; they do not yet confirm it.
\begin{itemize}[leftmargin=1.2em]
\item \textbf{Topological Data Analysis (TDA) in Neuroscience:} Researchers use TDA and algebraic topology to describe the shape of neural data and brain networks, including higher-order structures and cavities that pairwise graphs miss \cite{giusti,sizemore}. Project Ariel borrows this geometric spirit. Its UI could someday serve as an intuitive front-end for such ideas, but no study has yet linked on-screen line alignment to topological features of brain activity.
\item \textbf{Karl Friston's Active Inference:} Friston proposes that brains minimize free energy, an upper bound on ``surprise,'' by updating predictions (perception) or changing the world (action) \cite{friston}. When a user ``reanchors'' lines to make them parallel, the task resembles reducing a visible prediction error by action. This is an analogy for now, not a demonstration that the user is performing active inference.
\item \textbf{EMDR \& Spatial Working Memory Taxation:} Experiments indicate that taxing working memory while recalling a distressing memory (with eye movements or other demanding tasks) tends to reduce the memory's vividness and emotional intensity \cite{vdh}, and a meta-analysis found a moderate added effect of eye movements \cite{lee}. This evidence supports, but does not prove, the working-memory account. By replacing verbal therapy with the visual manipulation of lines and angles, the ``Reanchoring'' UI may act as an abstract form of working memory taxation, a hypothesis the pilot in Section 6b is designed to test.
\end{itemize}

\section{Postulates: Knowledge, Adaptation, and the Environment}
\label{sec:postulates}
\textit{In plain words.} Ryan's request was to treat a settled inner wisdom, even a wrong one, as the basis of a soul; to have the AI adapt without swallowing that wisdom as absolute truth; and to say that communication is a kind of existence, that choice begins as nonsense, and that a person must build an environment before communication can happen. This section states those ideas as postulates. It does not prove them. A metaphysical claim about the soul is not the sort of sentence a geometry app, or a citation, can prove. Saying that plainly is part of the respect owed to the claim.

\begin{itemize}[leftmargin=1.2em]
\item \textbf{P1 (faith and philosophy, not an empirical result).} Knowledge, in this framework, is formed as a person settles---fixes---their inner wisdom, even when that wisdom is wrong. That settled frame is the anchor of the self, and, in the author's Christian metaphysics, of the soul.
\item \textbf{P2 (a design rule).} The AI adapts to the person's frame and does not adopt their absolute truth. It is responsive. It is not sycophantic.
\item \textbf{P3 (a communication claim).} Communication is existence-in-relation. Choice begins as ``nonsense'' (not yet structured). Momentum turns that choice into structure that can be shared. The structure becomes an environment, and the environment shapes later choices.
\item \textbf{P4 (a communication claim).} A person has to create an environment---a shared context or frame---before they can communicate.
\end{itemize}

\textbf{Where the interface holds the postulates.} The Nonsense Layer is the origin: a locked picture of the thought model, including the anchors, $P_t$, $M$, $\gamma$, and entropy. It is the frame P1 points at, drawn so it can be seen and not edited by accident. Momentum on the status bar is the ``continuous momentum'' in P3. Freeform layers are the environment in P3 and P4: the person names a field and draws before the prompt (the act of communication) is built. The AI renames marks in that field's words and never moves a human shape, which is P2 in code. The canonical-terms toggle shows the geometry's own words, so the person's names are not treated as the last word.

\textbf{UI counterparts added in v1.1.} On the Nonsense Layer the person may \emph{establish} a self-point and a few anchors (values or beliefs), even if they might be wrong. The screen says ``Self established,'' with the date and what was fixed. The copy says establish, not prove. Only the person can write that record; the AI has no write path. It is included in the prompt and the frame export. That is P1 as a picture the person can point at, not a metaphysical demonstration. Each layer records who made the first move and who made the last. If the layer is blank, or the person has neither the first nor the last move, the meter reads AI-dominant. If the person takes the first or the last move, the person holds the environment and the AI only adapts. A small panel runs that comparison the same way every time (blank; AI moves; user first; user last) and shows the meter. The label says it is a model of P4 inside this app, not empirical proof. When the person's math line differs from the AI layer, the note says ``differs from.'' The AI keeps its own layer. It does not overwrite and it does not agree in silence. That is P2.

\textbf{P1: what the evidence can and cannot carry.}
Festinger described how people reduce the discomfort of holding two clashing cognitions, and how a choice can be followed by a firmer embrace of the chosen side \cite{festinger}. Ross, Lepper, and Hubbard showed that impressions of one's own ability can survive after the feedback that created them has been fully discredited \cite{ross}. Friston's free-energy account, and Clark's predictive-processing synthesis, describe brains as settling predictions by perception or by action \cite{friston,clark2016}. Together these support a narrow reading: people do lock in a view, including a wrong one, and minds can be described as settling. They do not show that the locked view is wisdom, and they do not show a metaphysical soul. P1's soul-language stays a confession of faith, kin to $V_s$, not a finding.

\textbf{P2: what the evidence can and cannot carry.}
Sharma, Tong, and colleagues reported that several assistant models, tuned with human feedback, often prefer answers that match a user's stated belief over answers that are simply true, and that both people and preference models sometimes reward that sycophancy \cite{sharma}. That is evidence for the risk P2 refuses. It is not evidence that this app has been measured against that risk. The engineering response is local and checkable: field words relabel; they do not rewrite geometry; the AI does not edit human layers; a toggle restores the canonical terms. Whether that is enough in a future model that talks in prose is an open question (below).

\textbf{P3: what the evidence can and cannot carry.}
Clark and Brennan argue that conversation is joint grounding: people coordinate content and process, and they cannot even start without some common ground and a medium \cite{clarkbrennan}. Vygotsky's account, in the Kozulin edition of \emph{Thought and Language}, treats speech as social in origin and only later internalized as verbal thought \cite{vygotsky}. Niche construction, as set out by Odling-Smee, Laland, and Feldman, is the biological claim that organisms change their surroundings and those changes feed back on later selection \cite{niche}. Read gently, these support three pieces: communication is relational work, structure can emerge from back-and-forth, and an environment can shape what comes after. They do not show that ``choice is the beginning of nonsense,'' and niche construction is not a study of this canvas. ``Nonsense'' here is our name for a frame that is not yet a message, not a clinical judgment.

\textbf{P4: what the evidence can and cannot carry.}
The same grounding chapter says that purpose and medium shape how common ground is built; without a setting, there is nothing to coordinate in \cite{clarkbrennan}. Gibson's affordances are what an environment offers an animal; perception, in that theory, is of those offerings, not of a bare sensation \cite{gibson}. Clark and Chalmers argue that some cognitive work sits outside the skull, in the tools and arrangements a person uses \cite{extended}. This supports P4 as a practical claim: a shared frame makes communication possible, and the frame can be part of the thinking. It does not prove a strict law that no signal can exist before an environment, and the extended-mind paper is philosophy, not an experiment on Ariel. In the app the practical reading is the one we build for: a layer (an environment) exists before a prompt is exported.

\section{Ten-Year Projection (2026--2036)}
\textit{In plain words.} The dates below are hopes, not forecasts. They stay only if later evidence earns them.
If Project Ariel is supported by evidence, its trajectory over the next decade could meaningfully shape mental health care. The phases below are aspirational goals, not predictions.
\begin{itemize}[leftmargin=1.2em]
\item \textbf{Phase 1: UI Prototyping \& TDA Validation (2026--2029).} Development of the 2D ``Reanchoring'' interface. Clinical trials utilizing Topological Data Analysis to test whether visual layer manipulation corresponds to reduced cognitive friction and free energy.
\item \textbf{Phase 2: Beyond the Categorical DSM (2030--2033).} A hoped-for transition from symptom-based diagnosis toward geometric state-space mapping, in the spirit of dimensional efforts such as NIMH's Research Domain Criteria \cite{rdoc}. The status quo baseline $M_{sq}$ is complemented by individualized topological mapping.
\item \textbf{Phase 3: Global System Floor Stabilization (2034--2036).} Wider adoption of Relativistic Cognitive Geometry. Cross-cultural mental health infrastructure informed by vector alignment, helping diverse populations achieve center-line balance without sacrificing their hidden cultural layers.
\end{itemize}

\section*{Eureka Log}
\textit{In plain words.} These are design insights from building the interface on 1 November 2026. They are not experimental results.

\begin{enumerate}[leftmargin=1.2em]
\item \textbf{The words belong at the edge.} A ribbon and a command line let the plane stay geometric. The research concepts (anchors, Euclid, patience) stay visible without writing a diagnosis on the drawing. Connects to the word-free Reanchoring protocol.
\item \textbf{Undo is a stack of whole states.} A stroke, an accepted suggestion, a merge, and a math line are edits. The partner's continuous motion is not. Because the AI never edits a human layer, undoing a human edit does not have to rewind the partner. Connects to preservation of layers: delete hides a shape into the stack; it does not pretend the stack was never there.
\item \textbf{One math line is the same drawing.} Points, segments, angles, parallels, circles, membership, intersection, and $I \to O$ already lived in the engine. Writing them as one Unicode sentence, and accepting that sentence back as an undoable edit, made the geometry readable without replacing it. Connects to the notation table.
\item \textbf{The Nonsense Layer is the origin, and it stays locked.} A reference frame that the person can see but not casually rewrite is what ``fixing'' a wisdom looks like in the tool, including when that frame is only a model. Other layers are free. Connects to P1 as a picture, not as a proof, and to P4: the frame exists before the other fields speak.
\item \textbf{A field changes the words, not the shapes.} Physics, music, psychology, theology, and math each have a vocabulary. The person's own names join it. The canonical toggle puts the geometry's words back. The partner explains; it does not agree that the person's names are the final truth, and it does not move their ink. Connects to P2.
\item \textbf{The prompt waits for an environment.} Export speaks only from a layer that already exists: notation, English, bag of words, math line, glossary, and alignment to the Nonsense Layer, in the order of the person's actions. Connects to P3 and P4.
\item \textbf{Establish, don't prove.} A self-point on the locked layer is something the person fixes in the tool, dated, and included when they speak. The AI cannot write it. The button says establish. Connects to P1 as a practice, not a proof.
\item \textbf{First or last move holds the room.} A blank layer reads AI-dominant. The person's first or last move flips the meter. A panel repeats the four cases so anyone can see the rule. If the pictures differ, the note says so and neither side is overwritten. Connects to P2 and P4 as a model inside the app.
\end{enumerate}

\section{Open Questions and Undefined Areas}
\textit{In plain words.} This list is an invitation. Undefined is not a failure. It is room for the next person, in a field we have not named yet.

\begin{itemize}[leftmargin=1.2em]
\item Does settling a wrong belief (P1's empirical edge) help, harm, or both, when the task is line alignment rather than a questionnaire? We do not know.
\item Does domain language reduce sycophancy, increase it, or leave it untouched once a model is allowed to answer in prose? The local app does not answer in prose. P2 is untested there.
\item Is ``nonsense'' a useful scientific word for an unstructured choice, or only a useful name inside this project? Open.
\item Which fields need terms we have not written? Law, ecology, nursing, mathematics education, and contemplative practice are untouched. A reader can add a map. The geometry does not have to change.
\item Do angle error, time to align, and path efficiency (Section 6b) track anything a person would call knowledge? That is a study, not a claim.
\item How should a locked reference layer be governed when two people share one canvas? Guardians are metaphors until a consent model exists.
\item $V_s$ is not a parameter to fit. How a study stays honest about optional faith, for participants who do not share it, remains an ethical design problem.
\end{itemize}

\section{Toward a New Field of Study}
\textit{In plain words.} If this work is a seed, the field is not ``AI therapy.'' It is the study of how a person and a tool settle a frame, keep the person's ink sovereign, and only then speak. Call it, provisionally, \emph{frame-first communication}. The name can change. The questions should be testable.

\begin{enumerate}[leftmargin=1.2em]
\item \textbf{Question.} When people align lines to a locked Nonsense Layer, do belief-perseverance measures (a debriefing task in the family of Ross, Lepper, and Hubbard) move differently than in a control that only watches lines? \textbf{Test.} A small pre-registered pilot, the shape of Section 6b, with SUDS and the in-app metrics, and with P1's soul-language left out of the participant materials unless a person opts in.
\item \textbf{Question.} If a prose assistant is given the exported prompt, does a canonical-terms instruction reduce agreement with a planted false belief, compared with a field-words instruction? \textbf{Test.} A vignette study modelled on the sycophancy tasks of Sharma and colleagues. The outcome is truthfulness against a planted error, not user satisfaction.
\item \textbf{Question.} Does requiring a named layer before export change the clarity of a later explanation, compared with exporting from a blank page? \textbf{Test.} Two arms, same geometry, communication scored by a rubric for common ground (Clark and Brennan's content and process), by raters who do not know the hypothesis.
\item \textbf{Question.} Can a new field (for example ecology or law) add a term map that a domain expert accepts, without any change to the geometry engine? \textbf{Test.} Expert review of the glossary against the unchanged math line. Success is ``the words fit and the picture is the same,'' not a clinical effect.
\end{enumerate}
None of these studies exist yet. Until they do, the postulates stay postulates. The app remains a model.

\section{Limitations and Ethics}
\begin{itemize}[leftmargin=1.2em]
\item \textbf{Conceptual status.} The master equation is a structured metaphor. Its terms lack units, measurement procedures, and parameter estimates, so it cannot yet be solved, fitted, or falsified as a whole. Only the UI-level proxies in Section 6b are testable today.
\item \textbf{Softened claims.} EMDR research supports working-memory taxation as one mechanism; it does not prove that any visual task heals trauma. TDA and active inference are compatible frameworks, not confirmations of Ariel. Replacing the DSM is a long-range aspiration; diagnostic systems change slowly and only with strong evidence.
\item \textbf{Symptoms as protection.} Reading isolation or self-pity as a protective ``Faraday cage'' can be compassionate, but it must never delay care for depression, anxiety, trauma, or suicidal thoughts.
\item \textbf{Word-free design.} Removing words reduces semantic bias but also removes nuance and the healing power of being heard. Reanchoring is meant to complement, not replace, conversation with trusted people and professionals.
\item \textbf{Faith and freedom.} $V_s$ expresses the author's Christian conviction. Participants in any study must be free to hold any belief; spiritual framing should be optional and never coercive.
\item \textbf{Privacy.} Session data are personal health information: local by default, minimal, and shared only with informed consent.
\item \textbf{Not a substitute for care.} This framework and any app built from it are not medical devices or treatment. If you are in crisis in the U.S., call or text \textbf{988} (Suicide \& Crisis Lifeline); in an emergency, call \textbf{911}.
\end{itemize}

\section{Conclusion: From ``For'' to ``With''}
Project Ariel begins with a hopeful claim: nothing in a person is wasted. Layers are hidden, not destroyed; defenses are guards, not failures; imbalance can be ridden like a spinning top rather than feared. The master equation gathers these intuitions into one picture, and the Reanchoring UI offers a simple, word-free way to practice them: take a tilted line and gently turn it until it runs parallel to something true.

The same movement sits at the heart of the paper's theology. Striving \emph{for} God adds momentum and friction, the $\gamma\,dE/dt$ of a soul trying to earn its place. Being \emph{with} God collapses the distance; the line is already parallel to the Universal Angle, and effort becomes presence. C.\,S.\ Lewis's \emph{The Screwtape Letters} \cite{lewis} imagines how easily good work can become a substitute for that presence, and Ecclesiastes reminds us that comparison and striving are \emph{hevel}, vapor \cite{eccl}. The science in this paper is still to be done, and it should be done carefully and honestly. But its posture is settled: the work is offered not to reach God, but from the rest He already gives.

\section*{Acknowledgments}
Grok Bot (AI assistant), 2026. Used for drafting, code, and literature search; not a primary source \cite{grok}. The claims, the postulates, and the theology remain the author's.

All glory be to Jesus Christ.

\begin{thebibliography}{99}\small
\bibitem{aaq} Bond, F.\,W., Hayes, S.\,C., Baer, R.\,A., Carpenter, K.\,M., Guenole, N., Orcutt, H.\,K., Waltz, T., \& Zettle, R.\,D. (2011). Preliminary psychometric properties of the Acceptance and Action Questionnaire--II. \emph{Behavior Therapy}, 42(4), 676--688. doi:10.1016/j.beth.2011.03.007
\bibitem{eccl} \emph{Ecclesiastes} 1:2. The Holy Bible (World English Bible, public domain).
\bibitem{friston} Friston, K. (2010). The free-energy principle: a unified brain theory? \emph{Nature Reviews Neuroscience}, 11(2), 127--138. doi:10.1038/nrn2787
\bibitem{giusti} Giusti, C., Ghrist, R., \& Bassett, D.\,S. (2016). Two's company, three (or more) is a simplex: Algebraic-topological tools for understanding higher-order structure in neural data. \emph{Journal of Computational Neuroscience}, 41(1), 1--14. doi:10.1007/s10827-016-0608-6
\bibitem{rdoc} Insel, T., Cuthbert, B., Garvey, M., Heinssen, R., Pine, D.\,S., Quinn, K., Sanislow, C., \& Wang, P. (2010). Research Domain Criteria (RDoC): Toward a new classification framework for research on mental disorders. \emph{American Journal of Psychiatry}, 167(7), 748--751. doi:10.1176/appi.ajp.2010.09091379
\bibitem{phq} Kroenke, K., Spitzer, R.\,L., \& Williams, J.\,B.\,W. (2001). The PHQ-9: Validity of a brief depression severity measure. \emph{Journal of General Internal Medicine}, 16(9), 606--613. doi:10.1046/j.1525-1497.2001.016009606.x
\bibitem{lee} Lee, C.\,W., \& Cuijpers, P. (2013). A meta-analysis of the contribution of eye movements in processing emotional memories. \emph{Journal of Behavior Therapy and Experimental Psychiatry}, 44(2), 231--239. doi:10.1016/j.jbtep.2012.11.001
\bibitem{lewis} Lewis, C.\,S. (1942). \emph{The Screwtape Letters}. London: Geoffrey Bles.
\bibitem{sizemore} Sizemore, A.\,E., Phillips-Cremins, J.\,E., Ghrist, R., \& Bassett, D.\,S. (2019). The importance of the whole: Topological data analysis for the network neuroscientist. \emph{Network Neuroscience}, 3(3), 656--673. doi:10.1162/netn\_a\_00073
\bibitem{gad} Spitzer, R.\,L., Kroenke, K., Williams, J.\,B.\,W., \& L\"owe, B. (2006). A brief measure for assessing generalized anxiety disorder: The GAD-7. \emph{Archives of Internal Medicine}, 166(10), 1092--1097. doi:10.1001/archinte.166.10.1092
\bibitem{vdh} van den Hout, M.\,A., \& Engelhard, I.\,M. (2012). How does EMDR work? \emph{Journal of Experimental Psychopathology}, 3(5), 724--738. doi:10.5127/jep.028212
\bibitem{wolpe} Wolpe, J. (1969). \emph{The Practice of Behavior Therapy}. New York: Pergamon Press.
\bibitem{clark2016} Clark, A. (2016). \emph{Surfing Uncertainty: Prediction, Action, and the Embodied Mind}. Oxford: Oxford University Press.
\bibitem{clarkbrennan} Clark, H.\,H., \& Brennan, S.\,E. (1991). Grounding in communication. In L.\,B. Resnick, J.\,M. Levine, \& S.\,D. Teasley (Eds.), \emph{Perspectives on Socially Shared Cognition} (pp.\ 127--149). Washington, DC: American Psychological Association. doi:10.1037/10096-006
\bibitem{extended} Clark, A., \& Chalmers, D.\,J. (1998). The extended mind. \emph{Analysis}, 58(1), 7--19. doi:10.1093/analys/58.1.7
\bibitem{festinger} Festinger, L. (1957). \emph{A Theory of Cognitive Dissonance}. Stanford, CA: Stanford University Press.
\bibitem{gibson} Gibson, J.\,J. (1979). \emph{The Ecological Approach to Visual Perception}. Boston: Houghton Mifflin.
\bibitem{grok} Grok Bot (AI assistant). (2026). Used for drafting, code, and literature search; not a primary source.
\bibitem{niche} Odling-Smee, F.\,J., Laland, K.\,N., \& Feldman, M.\,W. (2003). \emph{Niche Construction: The Neglected Process in Evolution}. Princeton, NJ: Princeton University Press.
\bibitem{ross} Ross, L., Lepper, M.\,R., \& Hubbard, M. (1975). Perseverance in self-perception and social perception: Biased attributional processes in the debriefing paradigm. \emph{Journal of Personality and Social Psychology}, 32(5), 880--892. doi:10.1037/0022-3514.32.5.880
\bibitem{sharma} Sharma, M., Tong, M., Korbak, T., Duvenaud, D., Askell, A., Bowman, S.\,R., et al.\ (2024). Towards understanding sycophancy in language models. \emph{International Conference on Learning Representations}. arXiv:2310.13548. \url{https://arxiv.org/abs/2310.13548}
\bibitem{vygotsky} Vygotsky, L.\,S. (1986). \emph{Thought and Language} (A. Kozulin, Ed.). Cambridge, MA: MIT Press. (Original work published 1934.)
\end{thebibliography}
\end{document}
